\documentclass[11pt]{article}
\usepackage[margin=25mm]{geometry}
\usepackage[T1]{fontenc}
\usepackage{lmodern,amsmath,amssymb,booktabs,array}
\usepackage[hidelinks]{hyperref}
\usepackage{microtype}
\emergencystretch=1em
\widowpenalty=10000
\clubpenalty=10000
\title{Self-Play Adaptation of Published Koi-Koi Transformers:\\
GPU Simulation and Paired Evaluation}
\author{}
\date{Working draft, 17 October 2026}
\begin{document}
\maketitle
\noindent\textbf{Draft status.} Development experiments and their independent
artifact audits are complete. The frozen final evaluation is running; this
version contains no final-test results or claim of final superiority. Author
information, venue formatting, final results, and isolated inference timing
remain to be completed before submission.

\begin{abstract}
We study whether self-play policy optimization can improve an existing Koi-Koi
Transformer without changing its deployment architecture or adding search.
Our implementation combines a C++/CUDA simulator with policy initialization
from published RL-WP weights, an explicit stochastic policy over the inherited
action scores, and a separate public-observation value network. A fixed
development study compares two temperature/learning-rate settings at two fresh
training seeds. After 20 critic-only and 500 policy updates, the selected
setting achieves draw-adjusted win scores of 51.805\% and 51.500\% against the
pinned assisted RL-WP comparator on 5,000 shared duplicate deals; both
pair-bootstrap 95\% interval lower bounds exceed 50\%. A colder policy with a
jointly reduced learning rate performs worse on both paired comparisons.
These are development findings, conditional on the evaluated models and deals.
We specify a separate frozen final test against both RL-WP and RL-Point.
A distinct twelve-round experiment completes 100,001,204 training matches;
that computation is not attributed to the selected eight-round models.
The study emphasizes rule compatibility, separation of training scale from
strength evidence, retained negative experiments, and auditable model selection.
\end{abstract}

\section{Introduction}
Koi-Koi combines hidden hands, a concealed stock, and repeated decisions about
whether to stop after scoring or continue for a larger gain. Match play adds
another tradeoff: an action that increases expected round points may reduce
the probability of finishing the match ahead. This makes the optimization
objective and the evaluation rules central to an empirical strength claim.

Guan et al. introduced phase-specific Transformer agents and released both
win-oriented RL-WP and point-oriented RL-Point policies
\cite{guan2023koikoi,pinnedkoikoi}. We treat the pinned implementation as a
concrete comparator. We do not infer a universal ranking from its reported
results or assume that all Koi-Koi implementations use compatible scoring.
Our question is narrower and measurable: can additional self-play improve
greedy, publicly assisted match performance against these frozen policies?

The method adapts the existing actors with clipped policy optimization
\cite{schulman2017ppo}. It uses an explicit temperature to convert their
scores into a sampling policy, and a freshly initialized critic for
player-relative match returns. The main empirical contribution under study is
an improvement at two separately trained model seeds, followed by an untouched
test of the fixed method. The systems contribution is an audited simulator
and experimental pipeline that distinguishes completed matches, round endings,
decisions, and search work. We also retain failures of earlier search and
training interventions rather than presenting the selected method as the
only experiment attempted.

\section{Rules, observations, and comparator}
\subsection{Versioned game protocols}
The primary protocol, \texttt{paper8}, follows the pinned public implementation:
at most eight rounds, 30 initial points per player, and termination when either
player's points reach zero or below. The separate \texttt{research12} protocol
uses twelve rounds and disables bankruptcy termination. Other round rules
are shared. Results from these protocols are not interchangeable.

A round deals eight cards to each hand, eight to the field, and leaves 24 in
the stock. A player discards, resolves a matching capture, reveals a stock
card, resolves that capture, and may be offered stop/continue after an
increase in scored combinations. Forced transitions are folded into the
simulator step. In the pinned scoring convention there is no seven-point
doubling or opponent-koi doubling. After an exhausted round without a scoring
stop, the initial dealer receives one point. The repository's versioned
\texttt{docs/RULES.md} supplies the complete scoring contract, including
nonstandard combinations inherited from the comparator.

The action space has 48 card actions and two stop/continue actions. Decision
phases are discard, capture from the hand play, capture from the stock reveal,
and koi. The two capture phases share an actor. Policies receive their own
hand and public state/history; the opponent's hand, stock order, random-number
state, and future events are excluded. A legal-action mask is applied before
sampling or greedy selection.

\subsection{Pinned opponents and public assistance}
The comparator revision is
\texttt{89d789565c7c36c84b74662e43001d960bf80901}. Original RL-WP and RL-Point
weights are loaded without further training. Source files and all six weight
files are hashed in the final protocol. The comparison uses the public
endgame guard corresponding to upstream assistance for both sides. This guard
restricts actions when a public match-outcome decision is solved, including
stopping with a winning final-round score. It does not inspect hidden cards.

The primary opponent is assisted RL-WP. Assisted RL-Point is a distinct
secondary opponent, because point-oriented and win-oriented training need not
produce the same ordering. All final match evaluations use strict CPU
inference and greedy actions. GPU feature/inference checks support training
compatibility but do not establish universal floating-point action identity.

\section{Self-play adaptation}
\subsection{Policy and value functions}
We initialize the three phase networks from the original RL-WP weights.
Their inherited raw action scores $q_\theta(o,a)$ are not probabilities.
For the publicly legal action set $\mathcal A(o)$, training samples from
\begin{equation}
 \pi_{\theta,T}(a\mid o)=
 \frac{\exp(q_\theta(o,a)/T)}
 {\sum_{b\in\mathcal A(o)}\exp(q_\theta(o,b)/T)},\qquad a\in\mathcal A(o).
\end{equation}
The public endgame guard can further narrow $\mathcal A(o)$. At evaluation,
the agent chooses a maximizing legal score, so a positive temperature does not
change the mathematical argmax for fixed scores. Temperature can nevertheless
change the learned policy through sampling and optimization. Dropout remains
disabled during both collection and gradient updates, preserving a consistent
policy for stored log probabilities.

The inherited actors retain the upstream 256-wide, two-layer Transformer
architecture. Native features reproduce the original 300-by-50 representation,
including its historical input conventions. A separate value network receives
a 560-dimensional public/own observation and uses two 256-unit tanh hidden
layers followed by a scalar output. Its output layer is initialized to zero;
it shares no actor parameters. The packed training observation has 15,561
entries: 560 ordinary features, one phase identifier, and 15,000 actor features.
Neither actor nor this critic receives privileged hidden-state inputs.

\subsection{Player-relative returns}
Both players use the same current stochastic policy during self-play. Rewards
are zero until match termination, when the acting player receives $+1$, $0$,
or $-1$ for a win, draw, or loss. We use an undiscounted, player-relative
adaptation of generalized advantage estimation \cite{schulman2015gae}.
Let $d_t$ indicate match termination and let $s_t$ be $+1$ when consecutive
decision states have the same acting player and $-1$ otherwise. Then
\begin{align}
 \delta_t &= r_t+(1-d_t)s_t V(o_{t+1})-V(o_t),\\
 A_t &= \delta_t+(1-d_t)s_t\lambda A_{t+1},\qquad \lambda=0.95.
\end{align}
Rollout truncation bootstraps from the next observation. Round boundaries
retain continuation value unless the match ends. A terminal transition masks
the value of the automatically reset next match. The sign follows the actual
next acting player, which need not alternate at every decision phase.

Advantages are normalized over each rollout batch before optimization.
We optimize the clipped PPO surrogate with clipping parameter 0.2, add
one-half the value mean-squared error, and subtract 0.02 times policy entropy.
The global trainable gradient norm is clipped at 0.5. Adam uses
$\epsilon=10^{-5}$ and separate actor/value learning-rate groups. The first
20 updates freeze actor parameters and fit only the critic; the next 500
updates train both. This warmup is part of the reported budget.

\subsection{Fixed development comparison}
Each update collects 128 decisions from each of 128 parallel environments and
uses three minibatch epochs with minibatch size 512. Each model therefore
receives exactly 8,519,680 simulator decisions in 520 updates. The value
learning rate is $10^{-4}$. The warm setting uses $T=0.05$ and actor learning
rate $3\times10^{-6}$; the cold setting uses $T=0.01$ and
$6\times10^{-7}$. Both settings are trained at seeds 6710001 and 6720001.
Only the final update-520 checkpoints enter the comparison. Intermediate
snapshots are retained for accounting and recovery, not strength selection.

This is a comparison of two parameter bundles, not an isolated temperature
ablation. Although the two settings have the same learning-rate/temperature
ratio, they can differ in sampled trajectories, entropy, local gradients, and
adaptive optimizer states. We do not attribute a difference solely to one
mechanism. Execution is eager float32 with TF32 disabled. The GPU is an
NVIDIA GeForce RTX 5070 Ti; the runtime records PyTorch 2.12.0+cu130 and NumPy
2.5.4. Evaluation is kept separate from training.

\section{Evaluation and selection}
\subsection{Duplicate deals and uncertainty}
For each seed we play two matches with candidate/opponent seats exchanged,
retaining the same seat-assigned deal streams and initial dealer. If the
candidate's two match scores are $w_{i0},w_{i1}\in\{0,1/2,1\}$, the observation
for inference is $x_i=(w_{i0}+w_{i1})/2$. The reported draw-adjusted win score
is $\bar x$. We resample whole pairs, not individual match legs, for 5,000
bootstrap replicates using random seed 99 and linearly interpolated 2.5th and
97.5th percentiles. Paired method comparisons resample the per-deal differences.
These intervals describe deal variation conditional on each fixed model;
they do not measure uncertainty over an arbitrary distribution of training
seeds or guarantee exact finite-sample coverage.

Point difference is final candidate points minus opponent points, averaged
across the two seats. Because total points remain 60, this is twice the
candidate's final-points-minus-initial-30 statistic used in the original paper.
Positive expected points are reported separately and cannot substitute for a
win-score criterion. Shared deals across models or opponents do not create
additional independent deal samples.

\subsection{Development and frozen final rules}
The development study uses 5,000 duplicate pairs with seeds
7,200,000--7,204,999 for the initial policy and four trained endpoints. An arm
qualifies only if both training seeds have an absolute win-score 95\% interval
lower bound above 0.5. A cold-over-warm benefit requires positive lower bounds
for both paired score differences. Selection chooses cold if it qualifies and
shows that replicated benefit; otherwise it chooses qualified warm, then
qualified cold, and otherwise no method. The canonical model always uses
training seed 6710001 of the selected arm. There is no best-seed ranking.

These rules were fixed before this study's outcomes, but the research program
had already explored other methods and development deals. Consequently,
development qualification is not untouched confirmation. In particular, an
earlier two-seed screen of the warm method at 1,000 pairs was inconclusive.
We retain that result and use fresh model seeds in the present comparison.

The frozen final test evaluates both selected models against assisted RL-WP
and, separately, assisted RL-Point. Each of four jobs plays 10,000 duplicate
pairs on seeds 100,000,000--100,009,999: 80,000 matches on 10,000 unique deal
streams. The primary criterion requires both RL-WP interval lower bounds to
exceed 0.5. The secondary RL-Point criterion is recorded separately. There is
no model reranking, training, extension, or early strength stopping after
accessing these final outcomes. Four two-thread CPU workers run concurrently;
their recorded times are not an isolated latency benchmark.

\section{Audited development results}
The initial policy obtains exactly 0.5 paired score and zero paired point
difference on all 5,000 development deals. This passes an outcome-compatibility
check on that corpus, not a proof that every action or trajectory is identical
to the comparator. All four training jobs and five evaluations finish before
the study-level decision is audited.

\begin{table}[ht]
\centering\small
\begin{tabular}{rlrrr}
\toprule
Training seed & Setting & Score (\%) & 95\% interval (\%) & Point diff. \\
\midrule
6710001 & warm & 51.805 & [51.090, 52.525] & 2.3802 \\
6710001 & cold & 50.755 & [50.135, 51.395] & 1.3046 \\
6720001 & warm & 51.500 & [50.765, 52.245] & 2.1546 \\
6720001 & cold & 50.565 & [49.930, 51.215] & 1.1602 \\
\bottomrule
\end{tabular}

\caption{Development results against assisted RL-WP. Every row uses the same
5,000 duplicate deal seeds. Point differences use candidate-minus-opponent units.}
\label{tab:development}
\end{table}

Both warm seeds qualify under the fixed development rule
(Table~\ref{tab:development}). Cold qualifies only at the first seed, and its
second interval includes 50\%. The paired differences favor warm on both
seeds (Table~\ref{tab:paired}); cold does not meet the replicated-benefit
criterion. The fixed rule therefore selects warm and the first-seed canonical
checkpoint. These observations support this bundle comparison at the measured
endpoints, not a general claim that reducing temperature harms self-play.

\begin{table}[ht]
\centering\small
\begin{tabular}{rrrr}
\toprule
Training seed & Cold minus warm (pp) & 95\% interval (pp) & Point diff. change \\
\midrule
6710001 & -1.050 & [-1.875, -0.200] & -1.0756 \\
6720001 & -0.935 & [-1.810, -0.075] & -0.9944 \\
\bottomrule
\end{tabular}

\caption{Paired development differences. Negative values favor warm.}
\label{tab:paired}
\end{table}

The earlier warm screen, using training seeds 6681001 and 6691001, produced
51.275\% [49.725, 52.850] and 51.075\% [49.475, 52.725] on 1,000 shared
development pairs. Neither cleared the original replicated criterion.
We do not pool these deals or models with the later study, extend those old
endpoints, or describe the present selection as a single unselected experiment.

\subsection{Final evaluation status}
At the time of this draft the final jobs are running. No final outcomes are
transcribed from live logs or incomplete reports. A completed independent
audit is required before populating the final results and revising the
abstract or conclusions. Failure of either model to clear the primary rule
will be retained as a failed replicated final criterion, regardless of point
gains or the secondary comparator's outcome.

\section{Systems, budgets, and negative evidence}
\subsection{Native simulation and validation}
Python reference rules and shared C++/CUDA transitions implement the same
versioned profiles. Differential tests exercise CPU/GPU transitions, public
observations, legal masks, and player-relative signs. Native public features
are checked against the pinned model inputs. A recorded compatibility corpus
contains 12,861 decisions from 100 complete random matches, 40 cached decision
states, and 128 hidden-state variants, totaling 13,029 active states. On this
finite corpus native CPU/CUDA features agree with the original inputs and
both original policies have zero observed assisted-action disagreements at
the tested batch shape. This is finite evidence, not a claim that numerical
near-ties can never differ elsewhere.

Checkpoints retain model weights, optimizer state, environment state, random
states, configuration, counters, and source identities. Resume checks cover
the exact next update on CPU and CUDA. Independent experiment auditors
recompute artifact digests, optimizer-step accounting, raw pair statistics,
and selection arithmetic; these audits do not independently replay every
training trajectory or every neural inference. The distinction matters when
assessing what a passed audit establishes.

\begin{table}[ht]
\centering\small
\begin{tabular}{rlrrr}
\toprule
Training seed & Setting & Matches & Round endings & Decisions \\
\midrule
6710001 & warm & 73,473 & 572,602 & 8,519,680 \\
6710001 & cold & 73,552 & 572,736 & 8,519,680 \\
6720001 & warm & 73,390 & 571,033 & 8,519,680 \\
6720001 & cold & 73,482 & 571,424 & 8,519,680 \\
\bottomrule
\end{tabular}

\caption{Actual terminal-event counts in the fixed development training.
Warmup is included; matches can cross the warmup/policy-update boundary.}
\label{tab:budget}
\end{table}

The four development jobs total 293,897 completed matches, 2,287,795 round
endings, and 34,078,720 decisions (Table~\ref{tab:budget}). Separately, a
\texttt{research12} baseline completes 100,001,204 new twelve-round matches,
1,200,025,665 round endings, and 18,157,928,448 decisions. That run satisfies a
computation-scale target, but it is not the selected pretrained policy's
training budget and does not prove public-baseline superiority. We do not
add shared initialization ancestry more than once or equate search rollouts
with training matches.

\subsection{Retained unsuccessful directions}
The broader development record includes architecture/critic variants,
koi-focused adaptation, direct published-opponent training, full-match return
experiments, and search. The repository's \texttt{docs/EXPERIMENTS.md} is the
chronological ledger; it contains pilot sizes, negative results, and subsequent
diagnostics. These studies informed method selection and must remain part of
the disclosure of exploratory work.

One fixed full-action rollout search extension improved the first tested model
but did not replicate on a second: its latter absolute score was 48.000\%
[46.225, 49.825], with a paired gain of 0.175 percentage points
[-0.750, 1.150] over no search. Independent confirmation samples for proposed
action changes also failed the both-model win-gain criterion. Separate
diagnostics found that selected-world gains shrink on independent samples and
that opponent-hand support assumptions can materially alter values at chosen
roots. Those selected-state diagnostics do not establish match-level benefits.
The final selected policy consequently uses no search.

\section{Limitations and interpretation}
Two trained seeds are evidence of replication at two models, not a reliable
estimate of performance across all training randomness. Model selection has
used substantial development experimentation. Final deal independence reduces
one source of selection bias, but does not erase that history or establish a
new algorithm's broad superiority. The bootstrap is conditional on each
evaluated policy and the stated deal generator.

The evaluated opponents and rules are deliberately specific. Beating the
pinned assisted RL-WP policy would meet the named baseline requirement; it
would not establish superiority over all public implementations, all scoring
variants, or humans. We have not measured exploitability or proved an
equilibrium guarantee. Public assistance is part of the evaluated policies,
and changing or removing it requires a separately labeled comparison.

The temperature study changes learning rate simultaneously and cannot isolate
temperature. GPU training and CPU evaluation use different arithmetic paths;
the compatibility corpus and recorded identities bound the claim rather
than eliminating all numerical uncertainty. The large twelve-round experiment
and the small pretrained eight-round adaptation answer different questions.
Finally, isolated decision latency is still pending. Concurrent final-test
throughput cannot be substituted for it.

\section{Current conclusion}
The audited development study selects warm pretrained-policy PPO at two fresh
model seeds, and the colder parameter bundle underperforms it in both paired
comparisons. The method is now frozen for an untouched two-comparator final
test. A final strength conclusion remains pending that test and its audit.

\appendix
\section{Reproduction and artifact identities}
The repository is \url{https://github.com/sheetau/koikoi}. The production package
source digest is
\begin{quote}\small\ttfamily
4b4360ed783c1d12cb5258e255f4b85f2e30dfa2e111fd8a078cf3439ce612cb
\end{quote}
The complete development audit digest is
\begin{quote}\small\ttfamily
6b8548605b50226738227ce83dba1a2e946efd282461b0fa8d48132a42cfe9ba
\end{quote}
The final protocol was frozen before launch with digest
\begin{quote}\small\ttfamily
4a69cff95079e2bd4cfcda84241f4096149c05cfe408bd11cbfde154eaa060e4.
\end{quote}
The warm checkpoints' full identities, all source inputs, exact unrounded
statistics, and separate scale evidence are in
\path{generated/evidence.json}. Tables are generated by running the module
\path{scripts.export_paper_evidence} with Python. The exporter verifies pinned
audit inputs and their report hashes without reading live final scores.
It transcribes evidence and is not itself an independent statistical audit.

The immutable study is run by \path{scripts.run_cold_ppo_screen} and checked
by \path{scripts.audit_cold_ppo_screen}. The final runner has separate
\texttt{--prepare} and \texttt{--run} modes, and refuses to replace an existing
run. The module \path{scripts.audit_final_published_test}, with the option
\texttt{--protocol-only}, checks the
freeze before launch; its full mode checks completed results. Never rerun an
existing final directory or change the protocol in response to final outcomes.
Reproduction on a separate machine should use a fresh checkout/workspace and
preserve exact dependency/model identities and explicit execution settings.

The Git repository currently excludes local checkpoints and run artifacts.
Consequently, code availability is not yet a complete archival data release.
A submission package still needs the permitted model artifacts, raw outcomes,
protocols, audits, dependency instructions, licenses, and attribution records,
with all identities checked against this ledger. No publication or external
data upload is implied by this draft.

\clearpage
\begingroup
\raggedright
\bibliographystyle{plain}
\bibliography{references}
\endgroup
\end{document}
